\section{What Was Hard}

\subsection{The Complexity Is Real}

The stack has many moving parts: deployment bundles, configuration overlays, custom resources,
object storage, ingress rules, secrets, and more. This is not a small installation, and there is
no point pretending otherwise. The complexity is the cost of having a platform that covers many
systems and can still be reviewed, versioned, and rolled back.

One lesson worth stating clearly: open source does not mean low effort. It means the effort moves
to the team. Someone has to understand upgrades, chart compatibility, storage sizing, label
conventions, retention policies, and user support. That is a good trade when the team needs control
and flexibility, but it should not be described as the easy path.

\subsection{Cardinality, Retention, and Cost}

Metrics and logs grow faster than expected. Labels are essential, but labels with too many unique
values make systems slow and expensive. Logs are valuable, but high ingestion rates with long
retention eat storage quickly. Managing this is ongoing work: setting ingestion limits, tuning
retention, sizing caches, watching query performance. A slow observability system is not useful
during an incident. The platform has to stay fast enough to be trusted when it matters most.

\subsection{Human Adoption}

Dashboards are only useful when people actually use them. That requires more than good tooling; it
requires naming things clearly, showing people where to look during real problems, and revisiting
dashboards that have stopped answering relevant questions. A platform that nobody looks at during
an incident is just a collection of graphs. The social side of adoption is as important as the
technical side.